\documentclass[11pt]{article}

\usepackage{amsmath,amssymb,amsthm}
\usepackage{graphicx}
\usepackage{hyperref}
\usepackage{geometry}
\usepackage{booktabs}
\usepackage{enumitem}

\geometry{margin=1in}

% Theorem-like environments
\newtheorem{definition}{Definition}
\newtheorem{proposition}{Proposition}
\newtheorem{constraint}{Constraint}

\title{\textbf{The Hilbert Substrate Framework III:}\\
No-Refolding and the Emergence of Three-Dimensional Space}

\author{Benjamin August Bray\\
\small Independent Researcher}

\date{January 2026}

\begin{document}

\maketitle

\begin{abstract}
We complete the Hilbert Substrate Framework by identifying a structural constraint---No-Refolding---that becomes unavoidable once stable matter exists within a Hilbert-space substrate governed by unitary dynamics, no-signaling, and no-forgetting. Building on earlier results demonstrating the emergence of spatial locality as an accessible attractor, we show that geometric dimensionality itself is selected dynamically through a settling process rather than global optimization.

Through numerical experiments on locality recovery across interaction graphs of varying intrinsic dimension, we find that three-dimensional geometry provides the deepest accessible basin compatible with stable fermionic matter. Crucially, the substrate does not optimize over dimensions; it settles at the first geometry capable of supporting persistent structure. Once matter forms, transitions to alternative geometries become dynamically inaccessible without violating information-theoretic consistency. Three-dimensional space therefore emerges not as a fine-tuned input, but as the first self-anchoring geometric configuration reachable under constrained unitary dynamics.
\end{abstract}

\section{Introduction}

\subsection{The Problem of Dimensionality}

Why does space have three dimensions? Standard physical theories typically treat dimensionality as a background assumption rather than a derived consequence. While anthropic arguments note that stable atoms and chemistry require three dimensions, they do not explain how such a dimensionality is selected dynamically. Approaches invoking extra compact dimensions similarly presuppose a higher-dimensional background without accounting for why three dimensions emerge as extended and stable.

The Hilbert Substrate Framework addresses this question from a different direction. Rather than assuming spacetime or dimensionality, it treats the universal Hilbert space and its unitary dynamics as fundamental, asking which forms of structure are unavoidable under minimal information-theoretic constraints.

\subsection{Context within the Hilbert Substrate Framework}

This work completes a three-part investigation:

\begin{itemize}
\item \textbf{Paper I} demonstrated that persistent heterogeneity in information propagation can arise dynamically under unitary evolution, providing an operational order parameter for emergent structure.
\item \textbf{Paper II} showed that spatial locality emerges as an accessible attractor on unitary orbits, even when the globally optimal (eigenbasis) description becomes kinetically inaccessible.
\item \textbf{This paper} addresses the remaining question: \emph{which} geometric structure is selected, and why dimensionality freezes at three.
\end{itemize}

\section{Constraints on the Hilbert Substrate}

We work within a minimal framework defined by three information-theoretic constraints. These are not dynamical laws but structural conditions required for observer-accessible physics.

\begin{constraint}[No-Signaling]
Local operations on one subsystem cannot be used to transmit information instantaneously to spacelike-separated subsystems.
\end{constraint}

\begin{constraint}[No-Forgetting]
Global evolution is unitary; information is never destroyed, only redistributed within the total Hilbert space.
\end{constraint}

\begin{constraint}[No-Refolding]
Once stable, structured matter exists, the accessible geometric configurations of the substrate are dynamically restricted. Existing structure cannot be arbitrarily refactored into alternative geometric arrangements.
\end{constraint}

Importantly, No-Refolding is not imposed independently. It emerges as a consequence of applying No-Signaling and No-Forgetting to substrates containing persistent structure. Any transition that would refold geometry while preserving matter would require either coordinated nonlocal operations or erasure of records, violating the first two constraints.

\section{Geometric Basins and the Settling Hypothesis}

\subsection{Accessibility versus Optimization}

As shown in Paper II, locality emerges because it is the simplest \emph{accessible} description of large quantum systems, not because it is globally optimal. The unitary orbit of a Hamiltonian contains many equivalent descriptions, but only a subset are dynamically reachable under locality-respecting transformations.

We extend this insight to dimensionality. The accessibility landscape contains multiple geometric basins corresponding to different effective dimensions. The substrate does not survey these basins globally; instead, it settles into the first basin capable of supporting persistent structure.

\subsection{Dimensional Requirements for Stable Matter}

Different spatial dimensions exhibit qualitatively distinct properties relevant to the stability of matter:

\begin{itemize}
\item \textbf{One dimension:} Exchange statistics are nonlocal, and stable bound structures are not supported.
\item \textbf{Two dimensions:} Braid statistics permit anyons, but stable atomic structure is precluded by both topological and dynamical considerations.
\item \textbf{Three dimensions:} Fermionic statistics emerge naturally, Coulombic bound states are stable, and topologically protected structures (e.g.\ knots) exist.
\item \textbf{Four or more dimensions:} While fermions are permitted, topological stability is weakened and matter fails to anchor geometry robustly.
\end{itemize}

\begin{proposition}
Three dimensions constitute the minimum spatial dimension supporting proper fermionic statistics, stable atomic bound states, and topologically persistent information storage.
\end{proposition}

\section{Numerical Evidence for Dimensional Settling}

\subsection{Experimental Design}

We perform numerical locality-recovery experiments on Hamiltonians defined over interaction graphs of varying intrinsic dimension. Each experiment consists of:

\begin{enumerate}[label=\arabic*.]
\item Constructing a local XX-model Hamiltonian on a specified graph.
\item Applying a global scrambling unitary.
\item Running a locality-biased recovery flow using graph-local two-qubit gates.
\item Measuring the reduction in nonlocal operator weight.
\end{enumerate}

\subsection{Results}

Across all tested system sizes and graph topologies, three-dimensional lattices consistently exhibit the deepest accessible locality basins. While higher-dimensional graphs are mathematically admissible, they are not dynamically accessed once stable structure forms at lower dimension.

These results do not rely on spectral dimension, which is poorly resolved at small system sizes. Instead, they reflect structural accessibility: how readily a given geometry traps the locality recovery flow.

\section{No-Refolding and Dimensional Lock-In}

Once matter forms within a three-dimensional geometric basin, further geometric transitions become dynamically inaccessible. Any attempt to refold the substrate into a different dimensional configuration would require either coordinated nonlocal transformations or destruction of existing records. Both are forbidden by the foundational constraints.

Dimensionality therefore freezes not because three dimensions are optimal, but because three dimensions are \emph{first}. The settling process terminates as soon as stable matter exists.

\section{Discussion}

\subsection{Relation to Other Approaches}

This framework differs fundamentally from anthropic or compactification-based explanations. Dimensionality is not selected by observers, nor by energetic preference, but by accessibility under constrained dynamics. Higher-dimensional geometries are not forbidden in principle; they are simply never reached.

\subsection{Limitations and Future Directions}

The numerical evidence presented here is limited to small system sizes. Extending these results to larger substrates will require tensor-network methods or analytical bounds on basin depth as a function of dimension. Future work may also explore cosmological implications of dimensional settling and potential signatures of early pre-geometric phases.

\section{Conclusion}

We have shown that three-dimensional space emerges naturally within the Hilbert Substrate Framework as the first geometry capable of supporting stable matter. Once such matter exists, No-Refolding locks the geometry in place. Dimensionality is therefore not an input to physics, but an emergent consequence of constrained unitary dynamics.

The universe does not choose three dimensions because they are optimal. It remains three-dimensional because this is where the settling process ends.

\bibliographystyle{plain}
\begin{thebibliography}{9}

\bibitem{HSF1}
B.~Bray.
\newblock \emph{The Hilbert Substrate Framework: Persistent Heterogeneity in Information Propagation}.
\newblock Zenodo, 2025.

\bibitem{HSF2}
B.~Bray.
\newblock \emph{The Hilbert Substrate Framework II: Accessibility and the Emergence of Spatial Locality}.
\newblock Zenodo, 2025.

\end{thebibliography}

\end{document}
